\section{Experimental Setup and Hyper-parameters}
\label{sec:Experimental-setup}
\subsection{Summarization Hyper-parameters}
To condense lengthy Indian Supreme Court judgments into structured and model-friendly inputs, we employed \texttt{Mixtral-8x7B-Instruct-v0.1}, a mixture-of-experts, instruction-tuned language model developed by Mistral AI. The summarization was conducted in a zero-shot setting using tailored legal prompts that extracted key elements such as facts, statutes, precedents, reasoning, and the final ruling.

The model was accessed via the HuggingFace Transformers interface and run on an NVIDIA A100 GPU with 80GB VRAM. Inputs were truncated to a maximum of 27,000 tokens to comply with the model’s context window. The output length was constrained to between 700 and 1,000 tokens to ensure consistency and legal completeness. A low decoding temperature of 0.2 was used to encourage determinism and factual alignment. These summaries served as inputs to the Retrieval-Augmented Generation (RAG) pipelines used for downstream judgment prediction and explanation.

\subsection{Judgment Prediction Hyper-parameters}
For the legal judgment prediction task, we used the \texttt{LLaMA 3–8B Instruct} model, which supports high-quality reasoning in instruction-following settings. The model was applied in a few-shot prompting setup without any task-specific fine-tuning. Input prompts consisted of structured summaries (produced by Mixtral) along with retrieved statutes and prior similar cases. These inputs followed a consistent legal instruction format to guide the model's prediction and explanation generation.

Inference was performed using the PyTorch backend with HuggingFace Transformers on an NVIDIA A100 GPU (80GB). The model was loaded using \texttt{device\_map=``auto''} for automatic device allocation. We used deterministic generation parameters (temperature = 0.2, top-p = 0.9) and controlled output format to ensure faithful and interpretable outputs. Each output consisted of a binary prediction (\texttt{0} for appeal rejected, \texttt{1} for appeal accepted/partially accepted) followed by a free-text legal explanation. No supervised fine-tuning was used, which allows our framework to be easily adapted to different legal datasets without retraining.

% \section{Inter-Annotator Agreement (IAA) for Expert Evaluation}
% \label{appendix:iaa_analysis}

% \subsection{IAA Metrics and Methodology}
% To ensure the reliability and consistency of expert ratings on the quality of generated legal explanations, we computed five widely accepted inter-rater agreement metrics:

% \begin{itemize}
%     \item \textbf{Fleiss' Kappa}~\cite{fleiss1971measuring}: Evaluates agreement among multiple raters for categorical judgments, adjusting for chance.
%     \item \textbf{Cohen's Kappa}~\cite{cohen1960coefficient}: Measures the pairwise agreement between two annotators, controlling for expected chance agreement.
%     \item \textbf{Intraclass Correlation Coefficient (ICC)}~\cite{shrout1979intraclass}: Assesses the degree of consistency in continuous ratings across multiple raters.
%     \item \textbf{Krippendorff's Alpha}~\cite{krippendorff2018content}: A versatile metric capable of handling varying scales and missing data, suitable for ordinal and interval data.
%     \item \textbf{Pearson Correlation Coefficient}~\cite{benesty2009pearson}: Quantifies the strength of the linear relationship between expert rating scores.
% \end{itemize}

% Three experienced legal experts independently rated a shared subset of model-generated legal explanations on a 10-point Likert scale, considering \textit{factual accuracy}, \textit{legal relevance}, and \textit{completeness}. The raters were blind to the model configurations to avoid bias and promote objective assessment.

% \subsection{IAA Findings and Observations}
% \label{appendix:iaa_observations}

% \paragraph{Interpretation:} 
% The overall inter-annotator agreement across all evaluation settings demonstrates moderate to substantial reliability. In the {Single} partition, pipelines such as \textit{CaseText + Statutes} and \textit{CaseText + Statutes + Precedents} achieved the highest agreement scores across most metrics (e.g., Fleiss' $\kappa >$  0.49, ICC almost 0.60), indicating stronger consensus among experts regarding their quality. This is aligned with the higher expert scores and other automatic evaluation metrics for these pipelines.

% In contrast, pipelines using only factual input or combining facts with retrieved statutes and precedents (\textit{CaseFacts Only} and \textit{Facts + Statutes + Precedents}) yielded relatively lower agreement scores (e.g., Fleiss' $\kappa < 0.35$, ICC $< 0.50$), reflecting the increased ambiguity or inconsistency in explanation quality when limited or noisy contextual information is used.

% The {Multi} partition exhibits slightly lower agreement metrics overall, potentially due to the complexity introduced by multiple judgment labels per case. Still, pipelines with richer legal context (\textit{CaseText + Statutes}, \textit{CaseText + Statutes + Precedents}) maintained comparatively higher consistency among annotators.

% \paragraph{Conclusion:} 
% These results reinforce the interpretability and credibility of our expert evaluation process. The observed agreement levels validate that the rating protocol is sufficiently robust to distinguish between explanation quality across different legal input pipelines. Moreover, the findings corroborate trends observed through both automatic and LLM-based evaluation metrics.

%%%%%%%%%%%%%%%%%%%%%%%%%%%%%%%%%%%%%%%%%%%%%%%%%%
% \begin{table*}[t]
% \centering
% \resizebox{0.9\textwidth}{!}{
% \begin{tabular}{lccccc}
% \toprule
% \textbf{Pipelines} & \textbf{Fleiss' $\kappa$} & \textbf{Cohen's $\kappa$} & \textbf{ICC} & \textbf{Kripp. $\alpha$} & \textbf{Pearson Corr.} \\
% \midrule

% \multicolumn{6}{c}{\textbf{Single Partition}} \\
% \midrule
% CaseText Only & 0.42 & 0.47 & 0.55 & 0.49 & 0.58 \\
% CaseText + Statutes & 0.51 & 0.55 & 0.61 & 0.57 & 0.65 \\
% CaseText + Precedents & 0.37 & 0.42 & 0.50 & 0.45 & 0.52 \\
% CaseText + Previous Similar Cases & 0.41 & 0.45 & 0.54 & 0.47 & 0.56 \\
% CaseText + Statutes + Precedents & 0.49 & 0.52 & 0.59 & 0.54 & 0.62 \\
% CaseFacts Only & 0.34 & 0.39 & 0.47 & 0.43 & 0.49 \\
% Facts + Statutes + Precedents & 0.29 & 0.34 & 0.42 & 0.38 & 0.44 \\

% \midrule
% \multicolumn{6}{c}{\textbf{Multi Partition}} \\
% \midrule
% CaseText Only & 0.40 & 0.45 & 0.53 & 0.48 & 0.57 \\
% CaseText + Statutes & 0.50 & 0.54 & 0.60 & 0.56 & 0.64 \\
% CaseText + Precedents & 0.36 & 0.41 & 0.49 & 0.44 & 0.51 \\
% CaseText + Previous Similar Cases & 0.39 & 0.44 & 0.52 & 0.46 & 0.55 \\
% CaseText + Statutes + Precedents & 0.47 & 0.51 & 0.58 & 0.52 & 0.60 \\
% CaseFacts Only & 0.32 & 0.37 & 0.46 & 0.41 & 0.48 \\
% Facts + Statutes + Precedents & 0.28 & 0.33 & 0.41 & 0.36 & 0.43 \\

% \bottomrule
% \end{tabular}
% }
% \caption{Inter-Annotator Agreement (IAA) statistics for expert evaluation of generated legal explanations across different pipeline settings for both Single and Multi partitions.}
% \label{tab:iaa_expert_scores}
% \end{table*}


%%%%%%%%%%%%%%%%%%%%%%%%%%%%%%%%%%
% hallucination_table
%%%%%%%%%%%%%%%%%%%%%%%%%%%%%%%%%%

\begin{table*}[ht]
\centering
\tiny
\resizebox{\textwidth}{!}{%
\begin{tabular}{|c|l|}
\hline

\textbf{Pipeline} &
  \multicolumn{1}{c|}{\textbf{Decision and Explanation}} \\ \hline
 
Ground Truth &
  { \begin{tabular}[c]{@{}l@{}}Decision: Rejected. \\Explanation:
  \\1. The parties to these appeals belong to Chokkura Thaliyadath tarwad, a hindu \\family governed by Marumakkathayam system of law as modified by the Madras \\Marumakkathayam Act,1932. They belong to the thavazhi of Cheriyammu Amma. \\Cheriyammu Amma acquired the suit property under a deed of gift Exhibit A-1 \\dated 19.6.1905.\\
  ........................................................\\
  Thus, in both the suits, the essential question that fell for decision, especially in the \\context of the only contention raised before us by learned Senior Counsel Mr. A.S. \\Nambiar, appearing on behalf of the appellant, was whether the transaction Exhibit-\\B-9 could be ignored by the sons of Kunhimalu Amma as a void transaction.", '6. \\Exhibit-B-9 is termed a deed of release\end{tabular}} \\ \hline

  
Case Text only &
  { \begin{tabular}[c]{@{}l@{}}Decision: Rejected.\\ Explanation:\\ The appeal is rejected because the lower courts have already upheld the decision that\\ the transaction Exhibit-B-9 is not void in law and the present court also dismissed\\ the second appeals filed by the plaintiff.\\
  ........................................................\\
  The plaintiff has failed to present any compelling argument or evidence that would\\ justify overturning the decision of the lower courts.  The appeal is therefore rejected.  \\The decision of the present court is well-supported by the law and the facts of the \\case.  The plaintiff has no valid grounds for appeal.\end{tabular}} \\ \hline


Case Text with Statutes &
  { \begin{tabular}[c]{@{}l@{}}Decision: Rejected \\ Explanation: 
  \\The appeal is rejected because the present court upholds the decision of the lower\\ courts, and the transaction Exhibit-B-9 is not void in law as it is a release of rights \\by all the other members of the thavazhi in favor of the only other member and \\not an assignment of undivided shares. \\
  ........................................................\\
  The plaintiff's argument that the transaction Exhibit-B-9 is void in law is not\\ supported by the evidence or the applicable laws, and the court's decision is a \\correct application of the law.  The court's decision is also consistent with \\the principles of the Marumakkathayam Law and the Madras Marumakkathayam \\Act, 1932.  The facts of the case do not support the plaintiff's argument that the\end{tabular}} \\ \hline

  
Case Text with Precedents &
  {\begin{tabular}[c]{@{}l@{}}Decision: Rejected \\Explanation:
  \\The appeal is rejected because the transaction Exhibit-B-9 is not void in law as it\\ is a release of rights by all the other members of the thavazhi in favor of the only\\ other member and not an assignment of undivided shares, which is in accordance \\with the Madras Marumakkathayam Act, 1932, and the ratio of the decision.\\
  ........................................................\\
  The ratio of the decision is that the transaction Exhibit-B-9 is not void in law as\\ it is a release of rights by all the other members of the thavazhi in favor of the\\ only other member and not an assignment of undivided shares.  The respondents \\did not have a subsisting right to be appointed on the basis of the select list prepared\\ on 4.4.87 as its life had expired. The Appointing Authority was justified in filling \\up the\end{tabular}} \\ \hline

  
  Case Text with Statutes \& Precedents &
  {\begin{tabular}[c]{@{}l@{}}Decision: Rejected \\Explanation:
  \\The appeal is accepted because the transaction Exhibit-B-9 is not void in law as it\\ is a release of rights by all the other members of the thavazhi in favor of the only\\ other member and not an assignment of undivided shares. \\
  ........................................................\\
  The Supreme Court's decision in this case is a significant step towards ensuring \\transparency and accountability in the electoral process, but it is not directly related \\to the issues in this case. Therefore, the appeal is accepted.  The decision in S. G. \\Jaisinghani vs Union Of India And Ors\end{tabular}} \\ \hline
\end{tabular}%
}
\caption{We observe `reject' judgments in the response generated by LLaMa-3x8B with varying supporting information.}
\label{cpt-hallucination-example}
\end{table*}
%%%%%%%%%%%%%%%%%%%%%%%%%%%%%%%%%%%%%%%%%%%%%%%%%%

% %%%%%%%%%%%%%%%%%%%%%%%%%%%%%%%%%%%%
%%%%%%%%%%%%%%%%%%%%%%%%%
% Indian Case Structure Example
%%%%%%%%%%%%%%%%%%%%%%%%%

% \section{Indian Case Structure Example}

\begin{table*}[t]
\centering
\begin{tabular}{|p{0.95\textwidth}|}
\hline
\textbf{CASE NO:} \\ \hline
 CIVIL APPEAL NO………………. OF 2014\\
(Arising out of SLP (C) No.22035 of 2013)\\ \hline
\textbf{APPELLANTS:} \\ \hline
Usha Bharti  \\ \hline
\textbf{RESPONDENT:} \\ \hline
State Of U.P. \& Ors \\ \hline
\textbf{DATE OF JUDGMENT:} \\ \hline
28/03/2014 \\ \hline
{ \textbf{BENCH:}} \\ \hline
{ Fakkir Mohamed Ibrahim Kalifulla} \\ \hline
\textbf{CASE TEXT:} \\ \hline

{ \begin{tabular}[c]{@{}l@{}}
... The earlier judgment of the High Court in the writ petition clearly merged with the judgment of\\the High Court dismissing the review petition. Therefore, it was necessary only, in the peculiar\\facts of this case, to challenge only the judgment of the High Court in the review petition. It.... \\ \\ 

...These Rules can be amended by the High Court or the Supreme Court but \textcolor{blue}{Section 114} can only \\be amended by the Parliament. He points out that \textcolor{blue}{Section 121 and 122}, which permits the High \\Court to make their own rules on theprocedure to be followed in the High Court as well as in...  \\ \\ 

...The principle of Ejusdem Generis should not be applied for interpreting these provisions. \\Learned senior counsel relied on Board of Cricket Control (supra). He relied on Paragraphs 89, \\90 and 91. learned senior counsel also relied on \textcolor{magenta}{S. Nagaraj \& Ors. Vs. State of Karnataka \& Anr}\\.[13] He submits finally that all these judgments show that justice is above all. Therefore, no...\\ \\ 


... We are unable to accept the submission of Mr. Bhushan that the provisions contained in Section \\28 of the Act cannot be sustained in the eyes of law as it fails to satisfy the twin test of reasonable \\classification and rational nexus with the object sought to be achieved. In support of this submission, \\Mr. Bhushan has relied on the judgment of this Court in \textcolor{magenta}{D.S. Nakara vs. Union of India}[16]. We...\end{tabular}} \\ \hline
\textbf{JUDGEMENT:} \\ \hline
{ \begin{tabular}[c]{@{}l@{}}.... When the order dated 19th February, 2013 was passed, the issue with regard to reservation was\\ also not canvassed. But now that the issue had been raised, we thought it appropriate to examine \\the issue to put an end to the litigation between the parties.
\\ \\
In view of the above, \textcolor{red}{the appeal is accordingly dismissed}.....\end{tabular}} \\ \hline
\end{tabular}%

\caption{Example of Indian Case Structure. Sections referenced are highlighted in blue, previous judgments cited are in magenta, and the final decision is indicated in red.}
\label{case-example}
\end{table*}
%%%%%%%%%%%%%%%%%%%%%%%%%%%%%%%%%%%%%%%%%%%%%%%


%%%%%%%%%%%%%%%%%%%%%%%%%%%%%%%%%%%%%%%%%%%%%%%
\begin{table*}[ht]
%\small
    \centering
    \begin{tabular}{|p{0.8\textwidth}|}
    \hline
{\bf Template 1 (prediction + explanation)}\\
\hline
   
{\bf prompt} = f``````Task: Your task is to evaluate whether the appeal should be accepted (1) or rejected (0) based on the case proceedings provided below..\\


{\bf Prediction}: You are a legal expert tasked with making a judgment about whether an appeal should be accepted or rejected based on the provided summary of the (case/facts) along with (Precedents/statutes/both) depending on the pipeline. Your task is to evaluate whether the appeal should be accepted (1) or rejected (0) based on the case proceedings provided below. \\


  {\bf case\_proceeding}: \# case\_proceeding example 1\\


  {\bf Prediction}: \# example 1 prediction \\


  {\bf Explanation}: \# example 1 explanation\\


  {\bf case\_proceeding}: \# case\_proceeding example 2\\


 {\bf  Prediction}: \# example 2 prediction\\


 {\bf Explanation}: \# example 2 explanation\\


{\bf Instructions}: L\#\#\# Now, evaluate the following case:
        
Case proceedings: {summarized\_text}
        
Provide your judgment by strictly following this format:

\#\#PREDICTION: [Insert your prediction here]\\
\#\#EXPLANATION: [Insert your reasoning here that led you to your prediction.]
           
Strictly do not include anything outside this format. Strictly follow the provided format. Do not generate placeholders like [Insert your prediction here]. Just provide the final judgment and explanation. Do not hallucinate/repeat the same sentence again and 
again''''''\\
\hline

    \end{tabular}
    \caption{Prompts for Judgment Prediction.}
    \label{tab:judgment_prediction_prompts_few}
\end{table*}
%%%%%%%%%%%%%%%%%%%%%%%%%%%%%%%%%%%%%%%%%%%%%%%


%%%%%%%%%%%%%%%%%%%%%%%%%%%%%%%%%%%%%%%%%%%%%%
\begin{table*}[ht]
\centering
\begin{tabular}{|l|}
\hline
\textbf{Instructions}:\\
You are an expert in legal text evaluation. You will be given:\\

A document description that specifies the intended content of a generated legal explanation.\\
An actual legal explanation that serves as the reference. A generated legal explanation that\\ needs to be evaluated.
Your task is to assess how well\ the generated explanation aligns with\\ the given description while using the actual document as a reference for correctness.\\
\\

\textbf{Evaluation Criteria (Unified Score: 1-10)}\\
Your evaluation should be based on the following factors:\\

\textit{Factual Accuracy (50\%)} – Does the generated document correctly represent the key legal\\ facts, reasoning, and outcomes from the original document, as expected from the description?\\
\textit{Completeness \& Coverage (30\%)} – Does it include all crucial legal arguments, case details,\\ and necessary context that the description implies?\\
\textit{Clarity \& Coherence (20\%)} – Is the document well-structured, logically presented,\\ and legally sound?\\
\\
\textbf{Scoring Scale:}\\
1-3 → Highly inaccurate, major omissions or distortions, poorly structured.\\
4-6 → Somewhat accurate but incomplete, missing key legal reasoning or context.\\
7-9 → Mostly accurate, well-structured, with minor omissions or inconsistencies.\\
10 → Fully aligned with the description, factually accurate, complete, and coherent.\\
\\
\textbf{Input Format:}\\
Document Description:\\
\{\{doc\_des\}\}\\
\\
\textbf{Original Legal Document (Reference):}\\
\{\{Actual\_Document\}\}\\
\\
\textbf{Generated Legal Document (To Be Evaluated):}\\
\{\{Generated\_Document\}\}\\
\\
\textbf{Output Format:}\\
Strictly provide only a single integer score (1-10) as the response,\\with no explanations, comments, or additional text.\\
\hline
\end{tabular}
\caption{The prompt is utilized to obtain scores from the G-Eval automatic evaluation methodology. We employed the GPT-4o-mini model to evaluate the quality of the generated text based on the provided prompt/input description, alongside the actual document as a reference.}
\label{tab:G-Eval Prompt}
\end{table*}

